\section{Paying the observable estimates on the same history}
\label{sec:common-observable-purchases}
The nonlinear return and the observable estimates have to be available
before a finite-volume history is selected. We now make the additional
choices in that order. The resulting statements are uniform over every
allowed representation on the purchased horizon. In particular, enlarging
the finite volume does not require a new choice of block lengths.

Retain the profiles, horizons and targets of
Section~\ref{sec:common-funded-history}. At a transition from stage $j$
to stage $i$, write $b_i$ for the outgoing support exponent, $e_{ji}$
for the support exponent of the local target itself, and $q_{ji}$ for
the rate of its input allowance. These are not the earned exponent
$b_i+3$ and the error exponent $-7/16$ requested of the output in
Section~\ref{sec:common-funded-history}: each target was purchased with
two positive support margins above its earned exponent and with a rate
strictly below its error exponent. The purchases therefore give
\[
             e_{ji}>b_i+3,\qquad q_{ji}\le-7/16.
\]
These are two different reserves: the first pays support enlargement,
whereas the second gives decay with the coupling. Fix a nondecreasing
integer occupancy schedule $N_j\ge1$. At each stage we need only the
two possible target stages and the finitely many grades through $j$.

\subsection{The broad parameter estimate and its localized price}
Put $y=x+s_G\log L$, $B=\lceil(\log x)^2\rceil$ and
$r=y^{1/256}$. Apply Section~\ref{sec:sector-analytic-prices} with
outgoing exponent $e=b_i+3$, the actual common support dictionary,
and the fixed derivative order and spatial power at stage $i$.
Its positive support reserve is $(e_{ji}-e)/4$.
The input allowance has upper exponent at most $-7/16$, so the
integrated constituent price has upper exponent $-2/5$ and the
rooted derivative price is eventually at most $y^{-3/8}$.
For each finite real test $\lambda$ on the parameter interval, its
ordinary logarithmic family $F$ therefore satisfies
\[
 \mathcal N_e\left(\int F(t)\,d\lambda(t)\right)
       \le y^{-3/8}M_{r/2}(\lambda).
\]
Here $\mathcal N_e$ is the rooted derivative budget at support
exponent $e$, and $M_{r/2}$ is the moment-test norm of
Section~\ref{sec:background-test-prices}. The forced family has
one common absolute coefficient majorant at radius $r/2$, with the
compulsory warning gain retained. The construction proves this for
all compatible input representations and all real parameter lines;
the majorant is not selected after the finite test.

The hypotheses on the real line deserve mention. Its coupling is the
same quadratic line used in Section~\ref{sec:real-coupling-interpolation}.
Its activity coordinates are affine, hence of degree at most two.
The path, cutoff, charts and input dictionaries are fixed on the line.
The enlarged real interval remains in the purchased band, by the
radius choice there. The compensated derivative and packet-action
estimates are supplied by Section~\ref{sec:parameter-comparison-suppliers}.
Thus all the local hypotheses of the sector estimate hold on these
actual lines, including their held fine fields and their cutoffs.

We check that the finite localization does not destroy the useful
power. Let $C_i$ be the signed-pullback derivative cost and define
\[
                  w(x)=C_i\{y^{-3/8}+4y^{-7/16}\}.
\]
This has upper exponent $-3/8$. The five auxiliary tori in
Section~\ref{sec:five-period-localization} have sides bounded by a
fixed multiple of $1+B$. Their numbers of sites, local ball counts,
and weighted changes of anchor are fixed powers of $1+B$.
The interpolation weights are fixed. For example, a contribution
from auxiliary torus $\Lambda_\nu$ has price at most
$2|\Lambda_\nu|w(x)$; its weighted change of anchor costs at most
$2^k(1+cB)^{pk}$ at derivative order $k$ and spatial power $p$.
The local allocation costs at most another fixed ball count and
$\exp((b_i+2)C_{\rm cover})$, with $C_{\rm cover}$ fixed.
Consequently both the localized raw derivative price and its
ordinary rooted price have upper exponent $-3/8$.
The latter is eventually at most $y^{-1/4}$, and the former at most
$y^{-1/8}$. These comparisons use strict exponent margins and are
uniform over the period and input representation.

The forced price survives the same localization with its forecast
normalization. The enlarged spatial factors have upper exponent zero.
The compulsory Gaussian warning is
$\exp(-\mu P(y)^2/c)$ with $\mu>0$ and fixed $c>0$; the permitted
forecast spends a strictly smaller part of that exponent. The ratio
of the spent price to the forecast tends to zero faster than any
fixed power of $y$, as in
Section~\ref{sec:nonlinear-payment-selection}. Thus the ordinary
price $y^{-1/8}$ and the soft correction can be paid simultaneously
in the forced absorption inequalities. This uses the chosen positive
forecast rate below its permitted fraction of $\mu$, already selected
before the delays. More precisely, the forecast hard rate is at most
$\mu/16$. Its soft allowance is twice the zero-order compulsory
activity price, $2A_{\rm comp}(x)$, and the resizing estimate is
applied at fraction $1/16$ before the final absorption. This keeps
both the factor two and the stated rate reserve. No new rate is
selected for an individual input.

\subsection{The fixed parameter radius and all derivatives}
Choose a fixed $q\ge4$ and a desired normalized output allowance
$0<s\le1/256$. The finite interpolation in
Section~\ref{sec:background-test-prices}, with degree $300$, has a
constant $U=U(300,q)$. Choose the small response parameter $t>0$
so that $6Ut\le s$. This choice precedes the nonlinear purchases;
Section~\ref{sec:common-funded-history} permits any positive such $t$.
The pointwise localization estimate is
\[
       3tC_{\rm st}x^{-11/10}+R(x),
       \qquad R\text{ has upper exponent }-23/20.
\]
Combining it with the just-proved broad bound gives the fixed-radius
price $sC_{\rm st}x^{-11/10}$. Indeed the discarded parameter tail
is at most
\[
        (1+U)(2q)^{301}y^{-365/256},
\]
which is $o(x^{-11/10})$, as is $UR(x)$. The linear reset and the
common forced coefficient majorant give the two normalized output
prices at radius $q/2$ from that section. The scalar projection has
price $sC_{\rm st}x^{-11/10}/\theta(x)^2$. Thus the finite-test
hypotheses of Section~\ref{sec:background-ray-estimates} are available
for the same return as its exact density identity.

The differentiated output also requires one common local collection
of constants before an input is supplied. This follows by making the
choices in Sections~\ref{sec:activity-jet-assembly} and
\ref{sec:parameter-comparison-suppliers} at the largest requested
order $k$. For the chart derivatives use a coarse radius
\[
 q_k(x)=Z_k\,2\theta_f(x),\qquad v(x)=(k+1)q_k(x),
\]
where $Z_k>0$ is the fixed angular chart constant and $\theta_f>0$
is the fine fluctuation threshold. Then $v(x)>0$, $kq_k(x)\le v(x)$,
and both radii have upper exponent $-1/2$. The comparison estimates allow any such
radius. Fix the cutoff convention before this collection of
constants is chosen. The first integrated polymer is the same
unmasked local function used in Section~\ref{sec:activity-jet-assembly}.
On a healthy support its mask is identically one in a neighborhood,
so its ordered derivatives are the derivatives of that unmasked
function. The angular-to-chart conversion at radius $q_k$ uses
$kq_k\le v$ and supplies all orders through $k$. This identifies
the pre-mask derivative estimate as well as the masked activity;
the former is required by later differentiated insertions. The Gaussian normalization estimates apply to every read
subset with the same row bounds, including the empty subset.
The enforcement and warning derivatives have respectively the
bounds proportional to the cover number and to one plus the
warning-support cover number; these are exactly the cutoff terms
retained in the differentiated activity construction.

Reserve a positive part of the target support exponent before
summing the differentiated factors. At every order through $k$,
its exponential weight absorbs the polynomial cover and length
factors. For completeness, choose a common positive support reserve
$\Delta=e_{ji}/2$ at outgoing exponent zero. If $A_k(x)\ge0$
is the common integrated derivative constant of upper exponent
zero and $u_k(x)>0$ the constituent allowance, its differentiated
allowance is
\[ s_k(x)=(1+2^kA_k(x))u_k(x). \]
It has the same negative upper exponent as $u_k$. The connected-tree
criterion $4(2B+1)^4s_k/\Delta\le1/2$ therefore holds eventually.
The ordinary and forced unweighted constituent allowances likewise
become at most $1/2$, and their tree criteria at reserve $\Delta$
hold at the same threshold. Since $0+\Delta+\Delta=e_{ji}$,
these are exactly the two reserves required in the renewed first
and second expansions. The unstrengthened construction in
Section~\ref{sec:high-degree-cluster-prices} proves their majorants
and all polynomial cover moments; we use it here at exponent zero,
with the same $\Delta$ in both expansions. Thus renewal does not
introduce an unrelated set of constituent prices.

The resulting connected activity majorants and nearby
integrable domination are those of
Sections~\ref{sec:activity-jet-assembly},
\ref{sec:gaussian-differentiation} and
\ref{sec:finite-volume-parameter-continuity}. Enlarge the finitely
many constants to common nonnegative ones and intersect their
large-coupling intervals. All those constants depend on the stage,
block length, chart and purchased prices, but not on the finite
volume or input representation. The subsequent unweighted total
sums may depend on the finite volume, as the continuity proof
allows. This distinguishes the uniform local derivative estimates
from the fixed-volume domination used to pass limits under sums.

\subsection{Insertion estimates and their common delays}
The primary, regular and forced one-insertion estimates use the
same local expansion. Their common heading and inverse bound were
proved in Section~\ref{sec:heading-history-handoff}. The primary
calibration uses the extra derivative available when $d+3\le i$;
the regular case uses its vanishing leading coefficient, and the
forced case retains its square-root forced-allowance normalization.
The auxiliary local integral and its comparison are those of
Sections~\ref{sec:principal-gaussian-comparison} and
\ref{sec:principal-comparison-jets}. The finite-support conditions
are imposed by increasing the shell and entry bounds, after their
fixed support radii have been selected.

For multiple insertions, the raised forced estimate and the
ordinary estimate have distinct prices. Use both, as in
Section~\ref{sec:finite-menu-history-handoff}. The tariff has
ordinary allowance $y^{-1}$, single-insertion allowance $y^{-2}$,
and forced allowance $2e^{-\mu P(y)^2/16}$. The occupancy and
projection inequalities of
Section~\ref{sec:projection-payment-producer} apply with $N=N_j$.
In particular the local projection bill is at most $e^{\sqrt y}$,
while its two error bills fit the fixed local splitting allowance.
The required geometric floor is imposed before $L$ is selected;
the occupancy and error inequalities are imposed afterwards.

Here is the common choice explicitly. At stage $j$ take the maximum
of the pre-length floors for the finitely many allowed target stages,
the primary comparisons, and the projected insertion geometry.
Choose $L_j$ above it. Having chosen wages, rates and protection,
choose the primary records with smallness at most $1/4$. For each
allowed target stage, take the finitely many eventual thresholds
for all grades through $j$, the broad and fixed-radius parameter
bounds, the differentiated output, the ordinary and raised menu
bounds, and the projection inequalities. Include all previous
thresholds. Replace their maximum at $j$ by its running maximum up
to $j$. Enlarge the shell bound by the required fixed comparison
radii and its growth lower bound by one.

Every conclusion used above is valid beyond its threshold, so all
are preserved by any further enlargement. The quantifiers are
therefore
\[
 \begin{gathered}
 \text{profiles and block floors};\quad\text{block lengths};\\
 \text{prices and protection};\quad\text{delays};\\
 \text{all compatible local inputs}.
 \end{gathered}
\]
The scale realization and the selected finite-volume history come
last. On that history each of the estimates applies to its actual
held transition; the endpoint representation, its density identity,
its singleton derivative and its marked profile all retain their
original common dictionary. In particular the observable estimates
do not choose another solution of the scalar shooting problem.

\subsection{Purchases for the terminal marked field}
Three further requirements enter the same order of choices. They are
consumed by Sections~\ref{sec:terminal-marked-construction},
\ref{sec:actual-marked-portfolio-calibration} and
\ref{sec:selected-centered-service}, and they are recorded with the
purchase, so that they hold on every schedule realized afterwards and
in particular on the selected run.

The first is the terminal fee. For each grade $d$ and each value of the
support exponent there is a threshold with the following property: at
every level $m$ of every schedule whose profile at $m$ has that support
exponent and whose coupling $\widehat\beta_m$ exceeds the threshold,
every torus side at level $m$ is at least the minimal side of the local
section of grade $d$, and the fee of that section at its camera constant
and at the forced allowance of level $m$ is at most the square root of
that allowance. This is the fee inequality of
Proposition~\ref{prop:terminal-marked-voice}; the threshold is produced
by the estimate at the end of
Section~\ref{sec:terminal-marked-construction}. At stage $j$ add the
nonnegative parts of these thresholds over all grades and all profile
indices through $j$, and include the sum $k_j$ in the entry delay of
stage $j$. Since the entry condition at level $m$ is imposed at the stage
$J(m+1)\ge J(m)$, every schedule paying the delays satisfies
$k_{J(m+1)}\le\widehat\beta_m$, and therefore the fee inequality at
every level $m$ for every grade $d\le J(m)$. The fee concerns the local
section of the family, with the family's label frames; it is fixed before
the block lengths.

The second and third concern the marked portfolio. Its sections and
marked families are fixed first, before the profiles' price maps, the
block floors and all later choices. For this fixed portfolio,
Section~\ref{sec:ex-marked-portfolio-calibration} produces a table of
nonnegative credits and stage thresholds $n_i$ on the block length, and
an entrance delay. Include the thresholds $n_i$ among the block floors
of stage $i$, together with the floors already listed; this is the
comparison of thresholds with floor lengths which that section requires
of its caller. Include its entrance delay among the entry delays, after
the protection parameters and before the final delay is selected. By the
argument given there, every schedule paying these delays then satisfies,
at every transition from stage $j$ to stage $i$ and every coupling $x$
beyond the entry delay of stage $j$, the profile-cost bound, the credit
bounds for all grades $d\le i$, and the tail inequality
\eqref{eq:ex-marked-height-jump}. The purchase records both facts
together with the portfolio: the credit table with thresholds not
exceeding the selected block floors, and the delay condition for the
selected delays. Neither requires a choice after the block lengths other
than the enlargement of the entry delays.

These three requirements are finite maxima of thresholds at each stage
and are stable under further enlargement, so they are compatible with
the common choice described above.

This supplies the additional local purchases needed to apply the
one-history transfer constructions. The passage from those local
purchases to the complete centered service and the ultraviolet
field family still requires their simultaneous global assembly.
